\documentclass[../../../include/open-logic-section]{subfiles}

\begin{document}
\olfileid[tr]{sth}{ordinals}{zfm}	
\olsection{$\ZFminus$: Bir Dönüm Noktası}

Yerine Koyma'nın nasıl gerekçelendirileceği (eğer gerekçelendirilebilirse)
kolay bir soru değildir. Bu yüzden soruyu \olref[replacement][]{chap}
bölümüne bırakacağız. Bununla birlikte Yerine Koyma'nın eklenmesiyle bir
başka önemli dönüm noktasına ulaştık. Artık $\ZFminus$ teorisi için gerekli
bütün aksiyomlara sahibiz. Ayrıntılı olarak:

\begin{defn}
$\ZFminus$ teorisinin aksiyomları Kaplamsallık, Birleşim, Çiftler, Kuvvet
Kümeleri, Sonsuzluk ve Ayırma ile Yerine Koyma şemalarının bütün
örnekleridir. Başka bir deyişle $\ZFminus$, $\Zminus$'a Yerine Koyma'yı
ekler.
\end{defn}
\noindent
Bu ad, \emph{Zermelo--Fraenkel} küme teorisini (daha sonra karşılaşacağımız
bir şey \emph{çıkarılmış} olarak) belirtir. Fraenkel bu onuru hak eder;
çünkü Yerine Koyma'nın \citeyear{Fraenkel1922} yılında formüle edilmesi ona
atfedilir, gerçi ilk kesin formülasyon \citet{Skolem1922} tarafından
verilmiştir.

\end{document}
